\section{Real-Data Assessment}
\label{sec:realdata}

The protocol above is simulation-based, and its central quantity---out-of-band energy---is defined against a known forward operator. Assessing how far the protocol extends to real measurements requires a system whose information support can be measured directly. We used a public experimental speckle dataset (measured through diffusers, paired ground-truth digits; 240 of its files sampled) and report what the protocol can and cannot deliver on it.

\paragraph{What is measurable.} The speckle power spectrum yields the system's information support: the grain full-width at half-maximum is $1.93$~px, corresponding to an MTF-support radius of roughly 52\% of Nyquist. The speckle intensity autocorrelation (Fig.~\ref{fig:robustness}, right) decomposes into a narrow grain peak (FWHM $\approx 1.9$~px) superposed on a wide platform extending to $\sim$$191$~px lag---the platform is the object information that speckle-correlation reconstruction classically exploits, and its stability across the dataset ($191 \pm 1.3$~px) identifies it as a system-level memory-effect/field-of-view quantity rather than an object property. A paired signal is detectable at the autocorrelation level: same-object matches separate from mismatched pairs (2D autocorrelation NCC $0.60$ versus $0.50$, Welch $t = 3.2$ over 15 pairs), so the dataset supports paired analysis. Pixel-level speckle--ground-truth correlations, by contrast, are too weakly separated to serve as a pairing statistic (means within ${\sim}0.01$ of each other over 20 pairs), which is why the paired analysis above operates at the correlation-structure level.

\paragraph{What is not yet available.} A quantitative bandwidth-protocol statement on this data---the real-data analogue of Sec.~\ref{sec:main-control}---requires a calibrated forward operator, and we did not obtain one under a minimal protocol. Three registration approaches failed systematically: (i) a cross-spectral transfer-function estimate degenerates, consistent with the recorded mapping not being a well-conditioned shift-invariant convolution at the object's recorded scale (the reliable structure lives in the autocorrelation domain, per the Siegert relation); (ii) radial autocorrelation matching is confounded because the system-level platform dominates the profile while the object's autocorrelation is not radially symmetric; (iii) two-dimensional autocorrelation scale matching produced a chaotic scale estimate (coefficient of variation $1.02$, values piling at both search boundaries). The failure pattern identifies the obstacle: this dataset was constructed so that the object-to-speckle mapping is learned end-to-end precisely because it is not analytically parsimonious, and our minimal registration protocol could not recover the operator from paired samples alone. Applying a bandwidth protocol to such data therefore requires a dedicated forward-calibration stage---either analytic or learned---which is a research task in its own right and is not a preprocessing detail. We report this boundary condition explicitly: the real-data section of this paper quantifies the environment (support, correlation structure, pairing statistics) and states plainly that the causal claim is established in simulation.
